\newacronym{fpga}{FPGA}{Field-Programmable Gate Array}
\newacronym{lqr}{LQR}{Linear Quadratic Regulator}

\chapter{GLOSSARY OF ACRONYMS}

\section{Building a Document Glossary}
  When using acronyms, you have to run LaTeX commands in a certain order to properly show them in the LIST OF ABBREVIATIONS glossary in the front matter.
  \begin{enumerate}
    \item Run pdflatex on your document to generate the .aux file.
    \item Run makeglossaries on your document to generate the .gls files.
    \item Run pdflatex on your document again to include the glossary in the front matter.
  \end{enumerate}
  If you set up \texttt{latexmk} to automatically build your document, it should run these commands in the correct order for you using the \texttt{.latexmkrc} file included in this example document.

  Keep in mind that you only need to run \texttt{makeglossaries} if you made some change to the list of acronyms, or if you reference a new acronym with the \verb|\gls| command for the first time.

\section{Subsequent References to Acronyms}
  The first time you refer to a \gls{fpga}, the full name is shown along with the acronym.
  Each subsequent usage just shows \gls{fpga}.

\subsection{Unused Acronyms}
  Unreferenced acronyms don't show up in the list of abbreviations in the front matter.

\section{Plural Acronyms}
  Using \verb|\glspl| instead of \verb|\gls| gives the plural form of the acronym.
  \glspl{lqr} is an example of a plural acronym.